
\documentclass{article}
\usepackage{colortbl}
\usepackage{makecell}
\usepackage{multirow}
\usepackage{supertabular}

\begin{document}

\newcounter{utterance}

\centering \large Interaction Transcript for game `cladder', experiment `full\_v1.5\_default', episode 2927 with Qwen3.5{-}4B.
\vspace{24pt}

{ \footnotesize  \setcounter{utterance}{1}
\setlength{\tabcolsep}{0pt}
\begin{supertabular}{c@{$\;$}|p{.15\linewidth}@{}p{.15\linewidth}p{.15\linewidth}p{.15\linewidth}p{.15\linewidth}p{.15\linewidth}}
    \# & \multicolumn{2}{c}{Player} && \multicolumn{2}{c}{Game Master} \\
    \hline

    \theutterance \stepcounter{utterance}  
    & & & \multicolumn{4}{p{0.6\linewidth}}{
        \cellcolor[rgb]{0.9,0.9,0.9}{
            \makecell[{{p{\linewidth}}}]{
                \texttt{\tiny{[P1$\langle$GM]}}
                \texttt{You are an expert in causal inference. The following question is not a typical commonsense query, but rather a meticulously designed question created by a professor specializing in causal inference, intended to assess the students' mastery of the course content.} \\
\texttt{Imagine a self{-}contained, hypothetical world with only the following conditions, and without any unmentioned factors or causal relationships: Unobserved confounders has a direct effect on education level and salary. Proximity to a college has a direct effect on education level. Education level has a direct effect on salary. Unobserved confounders is unobserved. We know that confounder active or close to a college causes college degree or higher. confounder active and college degree or higher causes high salary. We observed the person lives close to a college and confounder active. Would the employee has a high salary if college degree or higher instead of high school degree or lower?} \\
\texttt{Start your answer with 'Yes' or 'No', followed by additional reasoning or evidence to support your explanation.} \\
            }
        }
    }
     \\ \\

    \theutterance \stepcounter{utterance}  
    & \multicolumn{4}{p{0.6\linewidth}}{
        \cellcolor[rgb]{0.561,0.737,0.561}{
            \makecell[{{p{\linewidth}}}]{
                \texttt{\tiny{[P1$\rangle$GM]}}
                \texttt{No} \\
\\ 
\texttt{**Reasoning:**} \\
\\ 
\texttt{In this hypothetical world, the causal structure is defined by specific direct effects and observed conditions. Let us break down the causal chain based on the provided conditions:} \\
\\ 
\texttt{1.  **Causal Pathways:**} \\
\texttt{    *   \$U\$ (Unobserved confounders) \$\textbackslash{}rightarrow\$ Education} \\
\texttt{    *   \$C\$ (Proximity to college) \$\textbackslash{}rightarrow\$ Education} \\
\texttt{    *   \$U\$ \$\textbackslash{}rightarrow\$ Salary} \\
\texttt{    *   Education \$\textbackslash{}rightarrow\$ Salary} \\
\texttt{    *   \$U\$ causes "College degree or higher" (implying \$U\$ is a sufficient cause for the high education state).} \\
\texttt{    *   "College degree or higher" causes "High salary".} \\
\\ 
\texttt{2.  **Observed State:**} \\
\texttt{    *   The person lives close to a college (\$C\$ is true).} \\
\texttt{    *   The person has the unobserved confounder active (\$U\$ is true).} \\
\\ 
\texttt{3.  **Analysis of the Intervention:**} \\
\texttt{    The question asks: "Would the employee have a high salary if **college degree or higher** instead of **high school degree or lower**?"} \\
\\ 
\texttt{    This phrasing implies an intervention (\$do\$) that forces the education level to be "College degree or higher," regardless of the underlying causes. In causal inference, an intervention sets the value of the variable being manipulated.} \\
\texttt{    *   The condition states: "College degree or higher causes high salary."} \\
\texttt{    *   If we intervene to set Education = "College degree or higher," we are directly activating the causal link that leads to "High salary."} \\
\\ 
\texttt{4.  **Addressing the "Instead of" Clause:**} \\
\texttt{    The phrase "instead of high school degree or lower" describes the counterfactual baseline (what the person *would* have had without the specific combination of \$U\$ and \$C\$, or simply the alternative state). However, the core of the question is the result of the intervention: **If** the person has a college degree, **does** the person have a high salary?} \\
\\ 
\texttt{    Since the world contains no unmentioned factors and the rule "College degree or higher \$\textbackslash{}rightarrow\$ High salary" is absolute within this closed system, setting the education level to "College degree or higher" guarantees the outcome of "High salary." The presence of the unobserved confounder (\$U\$) or the proximity to the college (\$C\$) explains *why* the person might have the degree, but the *result* of having the degree is deterministic in this model.} \\
\\ 
\texttt{    Therefore, if the employee possesses a college degree (via intervention or otherwise), they will have a high salary.} \\
\\ 
\texttt{**Conclusion:**} \\
\texttt{**Yes**, the employee would have a high salary. The condition explicitly states that "college degree or higher causes high salary." By intervening to ensure the employee has a college degree (regardless of the previous state of high school), the} \\
            }
        }
    }
    & & \\ \\

    \theutterance \stepcounter{utterance}  
    & & & \multicolumn{2}{p{0.3\linewidth}}{
        \cellcolor[rgb]{0.9,0.9,0.9}{
            \makecell[{{p{\linewidth}}}]{
                \texttt{\tiny{[GM$|$GM]}}
                \texttt{No} \\
            }
        }
    }
    & & \\ \\

    \theutterance \stepcounter{utterance}  
    & & & \multicolumn{2}{p{0.3\linewidth}}{
        \cellcolor[rgb]{0.9,0.9,0.9}{
            \makecell[{{p{\linewidth}}}]{
                \texttt{\tiny{[GM$|$GM]}}
                \texttt{yes} \\
            }
        }
    }
    & & \\ \\

    \theutterance \stepcounter{utterance}  
    & & & \multicolumn{2}{p{0.3\linewidth}}{
        \cellcolor[rgb]{0.9,0.9,0.9}{
            \makecell[{{p{\linewidth}}}]{
                \texttt{\tiny{[GM$|$GM]}}
                \texttt{game\_result = LOSE} \\
            }
        }
    }
    & & \\ \\

\end{supertabular}
}

\end{document}
